\documentclass[12pt, a4paper]{article}

% Paquetes requeridos para el circuito
\usepackage[utf8]{inputenc}
\usepackage[T1]{fontenc}
\usepackage{circuitikz} 

\begin{document}

\begin{center}
    \begin{circuitikz}[american, scale=1.0, transform shape]
        
        % 1. Raíl inferior que representa el nodo común 'b'
        \draw (0,0) -- (8,0) node[circ] {} node[below] {b};
        
        % 2. Rama 1: V en serie con R1 entre 'b' y 'c'
        \draw (0,0) to[V, l=$V$] (0,2)
                    to[R, l=$R_1$] (0,4);
        
        % 3. Rama 2: Fuente de corriente I en paralelo (entre 'b' y 'c')
        \draw (2.5,0) to[I, l=$I$] (2.5,4);
        
        % 4. Rama 3: Resistencia R2 en paralelo (entre 'b' y 'c')
        \draw (5,0) to[R, l=$R_2$] (5,4);
        
        % 5. Raíl superior que conecta las ramas para formar el nodo 'c'
        \draw (0,4) -- (5,4) node[circ] {} node[above] {c};
        
        % 6. Resistencia R3 conectada entre el nodo 'c' y el nodo 'a'
        \draw (5,4) to[R, l=$R_3$] (8,4) node[circ] {} node[above] {a};
        
        % 7. Potenciómetro R conectado entre el nodo 'a' y el nodo 'b'
        \draw (8,4) to[vR, l=$R$] (8,0);
        
    \end{circuitikz}
\end{center}

\end{document}